Sia $\mathbb{X}$ un insieme ordinato. Prendiamo un suo sottoinsieme $\mathbb{E}$ con le seguenti
caratteristiche.

\[
    \begin{cases}
        \mathbb{E} \subseteq \mathbb{X} \\
        \mathbb{E} \neq \emptyset
    \end{cases}
\]

Posiamo quindi introdurre le seguenti definizioni per studiare gli estremi di $\mathbb{E}$.

\begin{definition}
    \textbf{Maggiorante per $\mathbb{E}$}

    Un elemento $M \in \mathbb{X}$ si dice maggiorante per $\mathbb{E}$ se e solo se $M$ è
    maggiore o al più uguale a ogni singolo elemento di $\mathbb{E}$.
\end{definition}


\begin{definition}
    \textbf{Minorante per $\mathbb{E}$}

    Un elemento $m \in \mathbb{X}$ si dice minorante per $\mathbb{E}$ se e solo se $m$ è
    minore o al più uguale a ogni singolo elemento di $\mathbb{E}$.
\end{definition}

\noindent\textbf{In simboli:}
\[
    \begin{aligned}
        M \text{ è maggiorante per } \mathbb{E} \Leftrightarrow
        \begin{cases}
            M \in \mathbb{X} \\
            \forall x \in \mathbb{E}, M \geq x
        \end{cases}
    \\ \\
        m \text{ è minorante per } \mathbb{E} \Leftrightarrow
        \begin{cases}
            m \in \mathbb{X} \\
            \forall x \in \mathbb{E}, m \leq x
        \end{cases}
    \end{aligned}
\]


\noindent\textbf{Conseguenze:}

Avremo quindi un insieme dei maggioranti

\[
    \magg{\mathbb{E}} := \left\{M\in\mathbb{X}\mid\forall x\in\mathbb{E}, M \geq x\right\},
\]
e un insieme dei minoranti
\[
    \mino{\mathbb{E}} := \left\{m\in\mathbb{X}\mid\forall x\in\mathbb{E}, m \leq x\right\}.
\]

\begin{observation}
    Gli insiemi $\magg{\mathbb{E}}$ e $\mino{\mathbb{E}}$
    possono anche risultare vuoti (non per forza contemporaneamente),
    nel momento in cui l'insieme $\mathbb{E}$ non ammetta maggioranti o non ammetta minoranti.
\end{observation}


\begin{definition}
    \textbf{Insieme superiormente limitato}
    
    L'insieme $\mathbb{E}$ si dice superiormente limitato se ammette almeno un maggiorante, ovvero se
    \[
        \magg{\mathbb{E}} \neq\emptyset.
    \]
\end{definition}

\begin{definition}
    \textbf{Insieme inferiormente limitato}
    
    L'insieme $\mathbb{E}$ si dice superiormente limitato se ammette almeno un minorante, ovvero se
    \[
        \mino{\mathbb{E}} \neq\emptyset.
    \]
\end{definition}

\begin{definition}
    \textbf{Estremo superiore}
    
    Dato un insieme superiormente limitato $\mathbb{E}$, si dice estremo superiore di $\mathbb{E}$ e
    si indica con $\sup \mathbb{E}$, il minimo, se esiste, dei maggioranti.
    \[
        \Lambda := \sup\mathbb{E}.
    \]
\end{definition}

\begin{definition}
    \textbf{Estremo inferiore}
    
    Dato un insieme inferiormente limitato $\mathbb{E}$, si dice estremo inferiore di $\mathbb{E}$ e
    si indica con $\inf \mathbb{E}$, il massimo, se esiste, dei minoranti.
    \[
        \lambda := \inf\mathbb{E}.
    \]
\end{definition}

\begin{definition}
    \textbf{Massimo di un insieme}
    
    Dato un insieme superiormente limitato $\mathbb{E}$ con estremo superiore $\Lambda := \sup\mathbb{E}$,
    si dice che $\Lambda$ è anche il massimo dell'insieme e si indica con $\Lambda = \max \mathbb{E}$ se
    $\Lambda \in \mathbb{E}$
    \[
        \lambda := \inf\mathbb{E}.
    \]
\end{definition}

\begin{definition}
    \textbf{Minimo di un insieme}
    
    Dato un insieme inferiormente limitato $\mathbb{E}$, si dice estremo inferiore di $\mathbb{E}$ e
    si indica con $\inf \mathbb{E}$, il massimo, se esiste, dei minoranti.
    \[
        \lambda := \inf\mathbb{E}.
    \]
\end{definition}

\begin{corollary}
    \label{cor:inf-mag}
    $M$ maggiorante per $\mathbb{A}$ $\implies$ $\forall M' > M$, $M'$ maggiorante per $\mathbb{A}$.
\end{corollary}
\begin{corollary}
    \label{cor:inf-min}
    $m$ minorante per $\mathbb{A}$ $\implies$ $\forall m' < m$, $m'$ minorante per $\mathbb{A}$.
\end{corollary}

Essendo che dato un elemento di $\mathbb{Q}$ ne esistono infiniti maggiori e infiniti minori, allora
per i corollari (\corref{cor:inf-mag} e \corref{cor:inf-min}) l'esistenza di un maggiorante $M$ per un sottoinsieme
di $\mathbb{Q}$ implica l'esistenza di infiniti maggioranti (maggiori di $M$) per quel sottoinsieme. Analogamente
l'esistenza di un minorante $m$ per un sottoinsieme di $\mathbb{Q}$ implica l'esistenza di infiniti minoranti
(minori di $m$) per quel sottoinsieme.



\begin{corollary}
    \label{cor:mag-inf-lim}
    $\magg{\mathbb{E}} \neq\emptyset \implies \magg{\mathbb{E}}$ inferiormente limitato.
\end{corollary}
\begin{demonstration}
    \textit{di} \corref{cor:mag-inf-lim}

    \[\Hp \magg{\mathbb{E}} \neq\emptyset \qquad \Th \magg{\mathbb{E}} \text{ è inferiormente limitato}\]

    \begin{list_2}
        \item Se $\magg{\mathbb{E}} \neq\emptyset$, segue locgicamente che
            \[ \mathbb{E}\subset \mino{\magg{\mathbb{E}}} \]
        \item Ed essendo necessario per l'esistenza di $\magg{\mathbb{E}}$ che $\mathbb{E} \neq\emptyset$, allora
            \[ \mino{\magg{\mathbb{E}}} \neq\emptyset\]
        \item Dunque $\magg{\mathbb{E}}$ è inferiormente limitato.
    \end{list_2}
    \qed
\end{demonstration}


\begin{corollary}
    \label{cor:min-sup-lim}
    $\mino{\mathbb{E}} \neq\emptyset \implies \mino{\mathbb{E}}$ superiormente limitato.
\end{corollary}
\begin{demonstration}
    \textit{di} \corref{cor:min-sup-lim}

    \[\Hp \mino{\mathbb{E}} \neq\emptyset \qquad \Th \mino{\mathbb{E}} \text{ è superiormente limitato}\]

    \begin{list_2}
        \item Se $\mino{\mathbb{E}} \neq\emptyset$, segue locgicamente che
            \[ \mathbb{E}\subset \magg{\mino{\mathbb{E}}} \]
        \item Ed essendo necessario per l'esistenza di $\mino{\mathbb{E}}$ che $\mathbb{E} \neq\emptyset$, allora
            \[ \magg{\mino{\mathbb{E}}} \neq\emptyset\]
        \item Dunque $\mino{\mathbb{E}}$ è superiormente limitato.
    \end{list_2}
    \qed
\end{demonstration}

\begin{example} \leavevmode

    \begin{list_3}
        \item Prendiamo come sottoinsieme di $\mathbb{Q}$
        \[ \mathbb{A} = \left\{ x \in \mathbb{Q} \mid 1 < x \leq 2 \right\} \]
        \item Essendo che $\forall x \in \mathbb{A}, 2 \geq x$, allora $2$ è un maggiorante per $\mathbb{A}$.
        \begin{enumerate}[label=\arabic*.]
            \item Dimostriamo che non esiste un maggiorante minore di 2
            \[ \Th \nexists M \in \mathbb{Q} \mid M < 2 \;\;\land\;\; \forall x \in \mathbb{A}, M \geq x .\]
            \item \label{pas:esi-mag/1} Supponiamo per assurdo che
            \[ \exists M \in \mathbb{Q} \mid M < 2 \;\;\land\;\;  \forall x \in \mathbb{A}, M \geq x .\]
            \item Allora, prendendo $r = 2$, si dimostra facilmente che
            \[ \exists r \in \mathbb{A} \mid M < 2,\]
            \[ \exists r \in \mathbb{A} \mid M < r.\]
            \item Viene contradetto il passaggio \pasref{pas:esi-mag/1}, l'$M$ così definito non può
            quindi essere un maggiorante per $\mathbb{A}$.
            \item Non esiste un maggiorante per $\mathbb{A}$ minore di $2$.
            \item $2$ è il minore dei maggioranti per $\mathbb{A}$.
            \qed
        \end{enumerate}
        \item Essendo $2$ il minore dei maggioranti per $\mathbb{A}$, allora
        \[ \sup\mathbb{A} = 2 \].
        \item Inolte $2 = \sup\mathbb{A} \in \mathbb{A}$, dunque
        \[ \max\mathbb{A} = \sup\mathbb{A} = 2 \].

        \item Essendo che $\forall x\in\mathbb{A}, 1 \leq x$, allora $1$ è un minorante per $\mathbb{A}$.
        \begin{enumerate}[label=\arabic*.]
            \item Dimostriamo che non esiste un minorante maggiore di 1
            \[ \Th \nexists m \in \mathbb{Q} \mid  m > 1 \;\;\land\;\; \forall x \in \mathbb{A}, m \leq 1 \]
            \item \label{pas:esi-min/1} Supponiamo per assurdo che
            \[ \exists m \in \mathbb{Q} \mid m > 1 \;\;\land\;\; \forall x \in \mathbb{A}, m \leq x .\]
            \item Allora, prendendo $r = 1$, si dimostra facilmente che
            \[ \exists r \in \mathbb{A} \mid m > 1,\]
            \[ \exists r \in \mathbb{A} \mid m > r.\]
            \item Viene contradetto il passaggio \pasref{pas:esi-min/1}, l'$m$ così definito non può
            quindi essere un minorante per $\mathbb{A}$.
            \item Non esiste un minorante per $\mathbb{A}$ maggiore di $1$.
            \item $1$ è il maggiore dei minoranti per $\mathbb{A}$.
            \qed
        \end{enumerate}
        \item Essendo $1$ il maggiore dei minoranti per $\mathbb{A}$, allora
        \[ \inf\mathbb{A} = 1 \].
        \item Inolte $1 = \inf\mathbb{A} \notin \mathbb{A}$, dunque
        \[ \nexists\min\mathbb{A}\].

        \item Rissumendo quindi i risultati abbiamo che
        \[
        \mathbb{A} = \left\{ x \in \mathbb{Q} \mid 1 < x \leq 2 \right\} \implies
        \begin{cases}
            \max\mathbb{A} = \sup\mathbb{A} = 2 \\
            \inf\mathbb{A} = 1 \\
            \nexists\min\mathbb{A}
        \end{cases}.
        \]
    \end{list_3}
\end{example}